% 第 1 章 作品功能与性能说明
% 本章引用键（实际使用）：fips203, kyber2018, guo2020, kyberslash, clangover, ghsa-f2v9, dudect, rtlf, ctgrind, microwalk, binsec
% 图表按章编号：本章为表 1.1–1.3
% 1.3 节的工作量总览为表 1.2，第 3 章的攻击可行性对照为表 3.1

\section{作品功能与性能说明}

\subsection{问题背景与赛题定位}

ML-KEM（模格密钥封装机制，Module-Lattice Key Encapsulation Mechanism）解封装的难点不在算法本身，而在源码级常数时间与产物级常数时间之间那条容易被忽略的错位。ML-KEM 由 CRYSTALS-Kyber 标准化而来，最终标准见 FIPS 203\cite{fips203,kyber2018}。解封装在 Fujisaki--Okamoto 变换下依次执行解密、重加密、比较与隐式拒绝，四步都接触秘密。任何一步出现依赖秘密的分支、访存或可变时延指令，处理耗时就会随秘密值变化，秘密信息便进入攻击者可以观测的时间通道。

这些问题有真实的案例。Guo 等对采用 FO 变换的密钥封装机制（KEM）给出密钥恢复的计时攻击，说明比较与拒绝这一步的时延可以携带信息\cite{guo2020}。KyberSlash 把秘密除法做成可复现的计时差异\cite{kyberslash}。clangover 表明编译器可以把源码里的无分支掩码还原成分支\cite{clangover}。

本作品围绕五类对象展开：liboqs 0.10.0 消息编码分支案例（M1）与 pq-crystals Kyber768 秘密除法案例（M2）对应真实历史实现；秘密下标查表人工校准夹具（M3）与提前返回比较人工校准夹具（M4）用于检验已知机制；mlkem-native 阴性对照（M5）用于观察误报。这五个对象承担三种互补角色：真实历史案例检验平台能否发现已有缺陷并核验官方补丁，人工夹具检验工具能否识别预先埋入的机制，阴性对照检查平台是否误报。

FO 变换把解密得到的消息重新加密，再与输入密文比较，比较结果决定输出真密钥还是假密钥。这条选择必须由常数时间比较与掩码完成，不能出现依赖秘密的分支。隐式拒绝的设计目的，是让失败路径与成功路径在外部不可区分。一旦比较或拒绝的路径依赖秘密，这个目的就落空：攻击者不需要私钥，只凭每次解封装的耗时，就能判断解封装走到了哪一步。

以 liboqs 0.10.0 的消息编码分支案例（M1）为例：未修复版在 57 µs 的完整解封装均值上多出 32.311 ns，95\% 置信区间 [20.282, 44.339]，只占均值的 0.06\%。% A-01 % B-01
这里比较的是全 0 消息与全 1 消息生成的两类合法等长密文，后文简称“主类对”。单次调用看不出这个差别——本机计时最小步进为 10 ns，32 ns 只相当于 3 个步进，一次调用的随机起伏就把它盖住了。把每类 10,000 次调用、三次独立运行的样本合起来，差异才从噪声里露出来。% A-30
因此，报告“未检出泄漏”时，还需要交代样本量和噪声水平，并给出最小可检测效应（MDE），即在指定检验水平和检出概率下，这套测量方案预计能分辨的最小时间差。它类似计时天平的分辨能力，不能被当作真实泄漏的上限。

C 语言源码没有写条件分支，并不代表编译后的机器码也没有分支。以 M1 为例，liboqs 0.10.0 的开发者用算术掩码代替条件分支，避免执行路径随秘密消息位变化。然而，在 x86\_64 上使用 clang 18 的 \texttt{-Os}（以缩小代码体积为目标的优化）编译时，这段掩码代码又被改写成了依赖消息位的条件跳转。该问题已登记为 CVE-2024-36405\cite{ghsa-f2v9}。相同源码在 ARM 64 位架构（aarch64）上生成的是位域提取和条件选择指令，没有出现这一秘密分支。% A-26 % A-27
因此，常数时间评估必须检查最终生成的机器码，并针对该编译产物进行测量。

换编译器、优化级别或目标架构都会改变产物，因此每次判定都绑定实际二进制与编译条件。

单工具下结论也不可靠。dudect 与 RTLF 只回答两类耗时是否可区分\cite{dudect,rtlf}，Memcheck 与 Microwalk 只回答哪里依赖秘密\cite{ctgrind,microwalk}，Binsec 只回答给定边界内是否常数时间\cite{binsec}，三类结论回答的问题不同，不能互相替代。第二个被测对象 M2（pq-crystals Kyber768 历史快照里的秘密除法）是最直接的例子：定位工具找到了三处操作数含秘密的整数除法，但完整解封装的合并耗时差为 $+4.58$ ns，95\% 置信区间为 [$-$3.08, 12.25] ns，跨过 0，表示数据尚不能将均值差与零区分。单独测量除法所在函数，合并差为 $+0.044$ ns，区间 [$-$0.430, 0.517] ns 也跨过 0。也就是说，代码里存在秘密相关指令，在当前机器和样本量下却未必能测出相应时间差。% A-03 % A-04
只用一个工具，M2 可能被读成存在泄漏，也可能被读成没有问题。同一份 M2，RTLF 在未修复版和已修复版均报告两类分布差异；dudect 的最大 $t$ 值落在已修复版，而采集前规定的均值检验在两个版本上均未检出显著差异。不同工具关注均值、分布形态或内部秘密依赖，结论需要结合各自检查的问题来解释。

还有一重困难在读者一侧。未检出常被当成常数时间，但阴性结论只在给定样本量、方差与检验水平下成立。不写检测预算，读者就不知道这个阴性排除了多大的效应。RTLF 在 42 个数据集上给出 30 个未检出、12 个检出\cite{rtlf}，% A-16
同一批数据换一种统计口径就可能换一种读法，光看“检出”还是“未检出”并不够。

M2 的检测预算为 24.7--67.1 ns，大于除法可能产生的效应。在这个预算下，未检出表示现有样本未提供足够的差异证据，不能说明除法本身没有时延。% A-15

这四类困难彼此牵连：效应相对总耗时越小，测量就越依赖严格控制与重复采样；源码层的写法只有经过实际编译产物核验，才能支撑产物层判断；不同工具给出的定位、统计与有界证明各自回答不同问题；检测预算若不公开，未检出结果便无法说明排除了多大的差异。因此，本作品构建检测、定位、防护与复测相衔接的工程流程，并将每项结论绑定到具体产物和检测预算。

题面 12 条把这条缝拆成可检查的动作：接入标准实现、公开向量验证、可重复计时、两类输入、统计检验、定位、防护、同环境复测、开销、环境实验、威胁模型与复现。任何一条只做一层，都会被后一条追问。本作品按检测、定位、防护、复测组织，让这些追问落在同一套证据上。

本作品以两类耗时是否可区分作为判定目标。每次判定绑定编译产物并按证据层登记，阴性结果同时报告该预算下的最小可检测效应。

\subsection{威胁模型与攻击面}

本作品评估的是本地白盒攻击者：有源码、能编译、能重复调用、能读完整解封装耗时，但看不到私钥、共享密钥与内部验证标志。表~\ref{tab:threat} 把能力假设与评估边界分成六项：攻击者能观测什么，判定止于哪一步。

\begin{table}[H]
\centering
\caption{威胁模型与评估边界}
\label{tab:threat}
{\small
\begin{tabularx}{\textwidth}{|>{\raggedright\arraybackslash}p{3.0cm}|>{\raggedright\arraybackslash}X|}
\hline
项目 & 本作品的设定 \\
\hline
评估者能力 & 持有源码与编译产物，可自选编译选项，可重复调用解封装，并读取每次调用的完整耗时 \\
\hline
秘密对象 & 解封装私钥中参与解密的部分、解密得到的消息、共享密钥、比较结果与隐式拒绝标志 \\
\hline
可观测量 & 每次完整解封装的耗时，两类合法等长密文的耗时分布 \\
\hline
不观测的量 & 私钥内容、共享密钥、内部验证标志，不使用功耗、电磁、故障注入，不假设与被测进程共享缓存 \\
\hline
攻击目标 & 判断两类密文的耗时分布是否可区分，并把差异归因到机制，不追求密钥恢复 \\
\hline
范围之外 & 真实远程链路、其他架构的计时层、功耗与电磁与故障注入、密钥恢复 \\
\hline
\end{tabularx}
}
\end{table}

表~\ref{tab:threat} 概括攻击者能力与作品边界。具体机制、所需观测和距密钥恢复的步骤见第 3.2 节的表~\ref{tab:attack}。

本作品在单台 x86\_64 裸金属、clang 18.1.3 上进行本地白盒评估，判断限定于指定产物与样本预算下的时间可区分性，不开展密钥恢复。M1 使用前期封存批次，环境与锁频规程相同，但不与 M2 至 M5 混合统计；aarch64 只用于机制层对照，不构成第二台测量机；网络实验是同机模拟，不代表真实远程。功耗、电磁、故障注入及跨环境合并均不在范围内。主类对、族外输入和攻击可行性分别见第 3.2 节与第 4 章。

\subsection{作品功能与系统组成}

平台围绕五个被测对象，打通“检测—定位—防护—复测”四个环节。统一计时、CPU 锁频、统计分析与定位结果登记四个公共模块，使每次运行都能追溯到具体对象、编译产物、测量环境和判定依据。

平台对外提供六项功能。首先是受控输入生成：按公开种子构造两类合法等长密文，使解密后的消息分别为全 0 和全 1，这两组主要比较对象称为“主类对”；篡改、格式破损和截断等扩展输入另行测试。其次是受控测量：固定 CPU 频率工作点，在预热后随机交错两类密文，逐次测量完整解封装并保存全部原始值。第三是统计检测：采集前写定输入类别、样本量和共同校正的检验集合（Holm 家族），再用 Welch 检验、Holm 误报校正和逆方差合并分析差异，同时报告 MDE。

另外三项是函数级定位、防护与审计、可复现。函数级定位用反汇编、秘密标记、可变时延检查、轨迹互信息与有界验证把阳性落到函数与指令。防护与审计覆盖七类手段，并在产物层核对目标机制是否真的消失。可复现把计划、产物摘要、原始表与重算脚本打包成一键入口。

五个被测对象分三类。M1 与 M2 是真实历史案例：前者是 liboqs 0.10.0 在 x86\_64 上被编译器引入的消息编码分支，后者是 pq-crystals Kyber768 历史快照里的秘密除法。M3 与 M4 是人工校准夹具，分别埋入秘密下标查表与提前返回比较，用来检验检测器能否识别已知机制。M5 是 mlkem-native 阴性对照，用来检查工具会不会误报。M1 对应 FIPS 203 初始公开草案（ipd），M2 是历史 Kyber768，两者都不等同于最终标准。五类对象的结果不能互相替代。

五类对象覆盖三种角色，是为了让同一套规程同时面对真实缺陷、已知机制与阴性对照。真实案例回答能否修复，人工夹具回答检测器能否识别，阴性对照回答会不会误报。三类问题不能用一个对象回答，也不能用一类对象的结论替另一类作保。

四个共享层是计时入口、锁频控制器、统计与功效、定位工具与结果登记。计时入口统一读钟与窗口，锁频控制器统一工作点与恢复核对，统计层统一家族、合并与预算口径，定位层统一工具版本与结果字段。四段链是检测、定位、防护、复测：检测回答能否分开，定位回答差异来自哪，防护回答改什么，复测回答改动是否生效。同一份原始数据可以供多层引用，结论只在其所属层成立。总体架构见图~\ref{fig:architecture}。

平台把判定用的二进制与计时用的二进制绑定为同一个。每次运行记录源摘要、产物摘要、编译器与选项，定位工具读的产物和计时窗口跑的产物是同一份。换编译选项等于换被测对象，必须重新登记，因为产物形态一变，基于旧产物的判定就不能直接沿用。这条约束把“源码改了”与“产物变了”分开，也是后面所有分层结论的前提。

工作量按对象、工具、排查、采集、样本、环境、测试、审计、统计与复现十个维度登记，汇总见表~\ref{tab:workload}。该表便于核查各对象的运行次数、逐格排查证据，以及采集记录能否回溯到具体运行。% A-22 % A-24 % A-31 % A-36 % A-37 % B-05 % B-06 % B-08 % B-09

\begin{table}[H]
\centering
\caption{工作量总览}
\label{tab:workload}
{\footnotesize
\begin{tabular}{|>{\raggedright\arraybackslash}p{2.6cm}|>{\raggedright\arraybackslash}p{4.2cm}|>{\raggedright\arraybackslash}p{8.4cm}|}
\hline
维度 & 数量 & 说明 \\
\hline
被测对象 & 5 个 & 2 个真实历史案例、2 个人工校准、1 个阴性对照 \\
\hline
定位工具 & 9 种 & 45 格矩阵中 31 格有真实运行产物 \\
\hline
全流程排查 & 90 格 & 3 实现 $\times$ 5 步骤 $\times$ 6 类，未检查 0 格 \\
\hline
正式采集 & 177 次控制器运行 & 54 次主实验 + 6 次 M2 热点微基准 + 117 次扩展对照 \\
\hline
主类对样本 & 每类 10,000 $\times$ 3 次 & 换种子、换物理核，原始值全量落盘 \\
\hline
优化级别 & 12 格 x86\_64 + 12 格 aarch64 & 3 种实现 $\times$ 4 级，产物判定与计时分开 \\
\hline
网络场景 & 5 档 $\times$ 3 种实现 $\times$ 2 次 & 本地直调、回环与低/中/高三档模拟 \\
\hline
功能测试 & 263 项 & 88 + 57 + 57 + 61 四项，只验证软件流程 \\
\hline
泄露渠道审计 & 128 文件 / 1,048 调用点 & 122 个接触秘密，0 条秘密泄露，1 条低危已修 \\
\hline
统计复核 & 4 个 Holm 家族 / 21 个成员 & 42 个 RTLF 数据集，每个阴性附 MDE \\
\hline
复现材料 & 精简与完整两档 & 匿名包 5/5，42 份原始长表，逐次运行台账 \\
\hline
\end{tabular}
}
\end{table}

\subsection{主要工作与创新点}

本作品把已有工具组织成以产物为中心、可审计、可分层的证据流程，不提出新的检测原语，也不把任何观察主张为新的缺陷。四项工作分别交代比较基线、具体改动、设计依据与实施代价，并共同约束结论的适用范围。

围绕上述四大挑战，本作品在证据建模、受控测量、多工具交叉定位与跨环境边界评估四个维度开展了系统性工作：

第一条是分层证据模型：把机制存在、热点可测、端到端可测与防护有效分为四层登记，并用最小可检测效应量化未检出。相对单工具、单层结论的评估方式，它要求同一个被计时的产物同时登记定位层、热点层、端到端层与复测层。每个阴性附预算，每个未执行单列。直觉容易把定位阳性当成计时阳性，M2 正好相反：三处除法定位阳性，端到端与热点微基准都未检出。四种状态是检出、未检出、未执行、不适用。这四种状态不是标签，而是决定一条结论能不能被引用的依据。代价是采集与工具运行增加，报告必须保留四种状态。证据见表~\ref{tab:modules}、表~\ref{tab:beforeafter} 与表~\ref{tab:mde}。

第二条是预注册与可审计的锁频协议：采集前冻结检验计划，控制器校验计划摘要并核对频率恢复，使每次判定可追溯、可复算。相对于事后选检验、只看 p 值的计时实验，主类对、样本预算、停止条件与 Holm 家族在采集前一次冻结。放行前做 20 次频率探针，退出核对恢复。锁频提供可核对、可恢复的频率工作点，减少动态频率变化对实验可比性的影响。三个对象的不锁频/锁频变异系数比分别为 0.964、1.006、1.089，说明本机的相对噪声水平接近，锁频的作用不能只用标准差是否下降来评价。% A-29
代价是机时与登记成本，见 3.5 节与表~\ref{tab:compare}。

第三条是多工具交叉与 90 格排查：用 9 工具 $\times$ 5 模块矩阵和 90 格排查，把定位阳性与计时阳性分开，并识别官方修复后残留的同源分支。相对单工具评估，它要求每个存在给出函数、指令与触发输入，每个不存在写明检查手段与覆盖输入。90 格中 5 格存在、75 格不存在、10 格不适用、0 格未检查。% A-22
90 格排查还发现，KyberSlash 官方修复提交后的产物仍保留与 clangover 同源的消息编码分支，说明只消除除法并不等于整个产物中已无秘密相关机制。该项是全流程排查方法的超额发现，尚未作逐指令时序归因，也不构成新的漏洞主张。代价是工具安装、版本固定与人工核对，见第 4.5、4.12 节。

第四条是产物与环境的依赖对照：把同一源码在优化级别、架构、虚拟化与网络条件下的产物形态与可检出性并列，界定结论的适用边界。相对单环境、单优化级别的评估，它把 12 格优化级别、12 格 aarch64、三台 WSL2（Windows 子系统 for Linux 2）与 5 档同机网络分开登记。时间证据不跨环境合并。直觉是分支越多越容易测到，实测 -O2 与 -O3 的命中位置更多，计时却未检出。M1 的机制在 aarch64 上不复现。% A-26 % A-27
代价是跨环境实验不能互相替代，见 4.10 节与 4.11 节。

四条工作合起来，回答的是同一个问题：同一份源码在不同产物上可以有不同结论，证据必须按产物、按层、按预算登记。每条结论都能回到一个具体的产物、一个明确的证据层和一组可复算的数字。这些结论不覆盖所有实现或所有环境。

\subsection{关键结果与性能概览}

五个模块在同一套测量与统计规程下形成了不同层次的结论：M1 的分支差异在修复后未检出，M2 虽在定位层命中但计时层未检出，M4 提供脆弱实现与常数时间实现的前后对照，M5 在既定预算内未误报。表~\ref{tab:modules} 展示五个模块在定位层、端到端计时层与防护复测层的横向结果。

五行共享同一计时入口、锁频工作点、预注册家族与功能门，因此测量规程可比；对象角色则各不相同，M1 与 M2 是历史缺陷，M3 与 M4 是人工夹具，M5 是对照。解读时应先比较各证据层，再结合对象角色理解每行结果，避免将规程上的可比性误作对象性质相同。

\begin{table}[H]
\centering
\caption{五模块分层结论}
\label{tab:modules}
{\footnotesize
\begin{tabularx}{\textwidth}{|>{\raggedright\arraybackslash}p{0.9cm}|>{\raggedright\arraybackslash}p{3.5cm}|>{\raggedright\arraybackslash}p{3.0cm}|>{\raggedright\arraybackslash}p{3.6cm}|>{\raggedright\arraybackslash}X|}
\hline
模块 & 对象与角色 & 定位层 & 完整解封装计时层 & 防护后 \\
\hline
M1 & liboqs 0.10.0 消息编码分支（FIPS 203 初始公开草案），真实历史案例 & 阳性：1 个秘密条件跳转 & 阳性：$+32.311$ ns [20.282, 44.339]，三次运行区间均不含 0 & 官方修复后分支消失，$-8.150$ ns [$-21.916$, 5.616]，未检出 \\
\hline
M2 & pq-crystals Kyber768 秘密除法（历史快照），真实历史案例 & 阳性：三处除法，1,280 次命中 / 3 处位置 & 未检出：$+4.58$ ns [$-3.08$, 12.25]，热点微基准 $+0.044$ ns [$-0.430$, 0.517] & 官方修复后除法计数归零，计时层前后均未检出 \\
\hline
M3 & 秘密下标查表，人工校准 & 阳性：1,280 次命中 / 2 处位置，互信息 1.922 bits & 未检出：$-0.73$ ns [$-5.56$, 4.10] & 消除查表版：0 报告，互信息 0.722 bits \\
\hline
M4 & 32 B 提前返回比较，人工校准 & 阳性：32 次命中 / 1 处位置，有界验证给出反例 & 阳性：脆弱实现 $-4.089$ ns [$-4.215$, $-3.963$] & 常数时间实现：$+0.055$ ns [$-0.363$, 0.473]，有界验证无秘密相关控制流 \\
\hline
M5 & mlkem-native，阴性对照 & 未命中：三种模式各 10 次，0 报告 & 未检出：$-5.30$ ns [$-11.77$, 1.16] & 无需修复，7 种工具误报 0 \\
\hline
\end{tabularx}
}
\end{table}

数字上的主线有四条。M1 未修复版合并差 $+32.311$ ns，区间 [20.282, 44.339]，三次运行区间均不含 0，官方修复后 $-8.150$ ns，区间 [$-21.916$, 5.616] 含 0。% A-01 % A-02
M2 完整解封装 $+4.58$ ns，区间含 0，除法热点微基准 $+0.044$ ns，区间含 0，定位层三处除法在修复后计数归零。% A-03 % A-04
M4 脆弱实现 $-4.089$ ns，区间 [$-4.215$, $-3.963$]，常数时间实现 $+0.055$ ns，区间含 0。% A-06 % A-07
M5 合并差 $-5.30$ ns，区间 [$-11.77$, 1.16] 含 0。% A-08

M2 按版本合并的均值差分别为 $+12.35$ ns 与 $-11.20$ ns，方向相反。衡量六次运行结果不一致程度的异质性指标 $I^2=58.3\%$，因此跨版本合并值只用于概括观测，正式判定仍采用采集前规定的逐次 Holm 检验，两版均未检出。% A-03
M4 的脆弱实现区间不含 0，常数时间实现区间含 0，前后对照因此在同一张表上完成。% A-06 % A-07
每个数字背后是三次独立运行与每类 10,000 次调用。三次运行的区间列在附录 D，正文只给合并值与区间。

未检出的格子都附了预算：M2 的 MDE 为 24.7--67.1 ns，M3 为 17.7--33.7 ns，M5 为 13.5--28.3 ns。% A-15
五个模块的结论按证据层分别解读：M1 与 M4 在定位层和计时层都发生变化，构成前后对照；M2 虽在定位层命中，计时层修复前后均未检出；M3 的开启查表版与消除查表版用于验证检测器能否识别已知机制；M5 是阴性对照。M4 是唯一在定位层与计时层均阳性、且常数时间实现均未检出的正对照，可据此检查检测链路对已知机制的响应。

M1 的三次运行区间均不含 0，Holm 家族中只有这一项被拒绝，这是本作品唯一在端到端层被判阳性的真实案例。% A-09

这些结论合起来说明一件事：修了与测不出是两件事。M1 的官方修复使分支消失，计时层随之回到区间含 0。M2 的除法在定位层归零，计时层却从未给出可区分的基线。M5 没有任何阳性，也只能说明本预算下未检出。三句话分开，读者才能判断每一步证据的强度。

表~\ref{tab:modules} 展示五个模块在定位层、端到端计时层与防护复测层的横向对比结果；各层分别记录机制证据、预注册统计判定和修复后变化。

性能开销按成对分块测。M1 修复使完整解封装 $+1.51\%$，M2 修复 $-2.92\%$：修复并不总意味着变慢，更快也不等于更安全。% A-25
M3 的消除查表版相对开启查表版为 $+3.60\%$，区间 [0.36, 9.98] 较宽，按描述性结果报告。M4 是 32 B 比较调用，首字节与末字节的相对开销分别为 $+26.21\%$ 与 $+6.85\%$；M1、M2、M3 的开销来自完整解封装，M4 来自 32 B 比较调用，均不含键生成与封装。

本章介绍作品入口，方法见第 3 章，逐项结果见第 4 章。赛题 12 项要求均有对应实现和证据，并在附录 A 逐项登记为满足。
